\begin{corollary}[7.1.5]
\label{IV.2.s7.1.5-fr}
Sea $A$ un anillo local noetheriano tal que exista
un $A$-módulo de tipo finito $M$ que sea un
$A$-módulo de Cohen-Macaulay y tal que
$\operatorname{Supp}(M)=\operatorname{Spec}(A)$
\textup{(como ocurre, por ejemplo, si $A$ es un anillo de
Cohen-Macaulay)}. Entonces $A$ es formalmente
equidimensional. Por tanto, según \textup{(7.1.4)}, para
que un anillo cociente $B$ de $A$ sea formalmente
equidimensional es necesario y suficiente que sea
equidimensional.
\end{corollary}

\begin{proof}
\oldpage[IV-2]{185}
En efecto, $\widehat M=M\otimes_A\widehat A$ es un
$\widehat A$-módulo de Cohen-Macaulay (\textbf 0, 16.5.2)
cuyo soporte es $\operatorname{Spec}(\widehat A)$;
por consiguiente, $\widehat A$ es equidimensional
(\textbf 0, 16.5.4).
\end{proof}
